\documentclass{letter}
\usepackage{eurosym}
\usepackage{fancyhdr}
\usepackage{graphicx}
\usepackage[utf8]{inputenc}
\usepackage{lmodern}
\usepackage[colorlinks]{hyperref}
\usepackage{eurosym}
\usepackage{longtable}
\usepackage[table]{xcolor}


\addtolength{\voffset}{-0.5in}
\addtolength{\hoffset}{-0.3in}
\addtolength{\textheight}{2cm}

\pagestyle{fancy}
\setlength{\headheight}{107pt}
%Variabili per ruoli
\newcommand{\ana}{\emph{Analista}}
\newcommand{\anas}{\emph{Analisti}}
\newcommand{\progs}{\emph{Progettisti}}
\newcommand{\RdP}{\emph{Responsabile di Progetto}}
\newcommand{\progr}{\emph{Programmatore}}
\newcommand{\progrs}{\emph{Programmatori}}
\newcommand{\ver}{\emph{Verificatore}}
\newcommand{\vers}{\emph{Verificatori}}
\newcommand{\gruppo}{\emph{Sons Of SWE}}
%variabili per documenti
\newcommand{\AdR}{\emph{AnalisiDeiRequisiti\_v4.0.0}}
\newcommand{\SdF}{\emph{StudioDiFattibilità\_v1.0.0}}
\newcommand{\VI}{\emph{VerbaleInterno}}
\newcommand{\NdP}{\emph{NormeDiProgetto\_v5.0.0}}
\newcommand{\PdP}{\emph{PianoDiProgetto\_v5.0.0}}
\newcommand{\PdPv}{\emph{PianoDiProgetto\_v5.0.0}}
\newcommand{\PdQ}{\emph{PianoDiQualifica\_v5.0.0}}
\newcommand{\PdQv}{\emph{PianoDiQualifica\_v5.0.0}}
\newcommand{\G}{\emph{Glossario\_v3.0.0}}
\newcommand{\MU}{\emph{usermanual\_v2.0.0}}
\newcommand{\MM}{\emph{ManualeManutentore\_v1.0.0}}
\newcommand{\VE}{\emph{VerbaleEsterno}}
\newcommand{\MS}{\emph{developermanual\_v2.0.0}}

\newcommand{\VIuno}{\emph{Verbale\_I\_1\_2018.03.05}}
\newcommand{\VEuno}{\emph{Verbale\_E\_1\_2018.03.07}}
\newcommand{\VIdue}{\emph{Verbale\_I\_2\_2018.03.08}}
\newcommand{\VEdue}{\emph{Verbale\_E\_2\_2018.03.28}}
\newcommand{\VItre}{\emph{Verbale\_I\_3\_2018.04.05}}
\newcommand{\VEtre}{\emph{Verbale\_E\_3\_2018.04.28}}
\newcommand{\VIquattro}{\emph{verbale\_I\_4_2018.04.30}}
\newcommand{\VEquattro}{\emph{verbale\_E\_4\_2018.07.05}}
\newcommand{\VIcinque}{\emph{verbale\_I\_5\_2018.05.22}}
\newcommand{\VIsei}{\emph{Verbale\_I\_6\_2018.06.20}}



%Variabili generiche
\newcommand{\indirizzo}{\emph{nome@cognome.it}}
\newcommand{\emailgruppo}{sonsofswe.swe@gmail.com}

\usepackage{xpatch}
\xpatchcmd{\makeletterclosing}{[3em]}{[1em]}{}{}
\patchcmd{\makeletterclosing}{[3em]}{[1em]}{}{}

\fancypagestyle{firstpage}{
	\fancyhf{}
	\renewcommand{\headrulewidth}{0.0pt}
	\renewcommand{\footrulewidth}{0.3pt}

	\fancyhead[L]{
		\includegraphics[height=3cm, keepaspectratio=true]{../template/img/logo.png}
	}
	\fancyhead[R]{
		\Large{\textbf{Gruppo \\ \textbf{\gruppo{}}}} \\
		\vspace{55pt}
	}

	\cfoot{ 
		\small{
			Gruppo: \gruppo{} \\ e-mail: \emailgruppo{}
		}
	}
}
%%%%%%%%%%%%%%%%%%%%%%%%%%%%%%%%%%%%%%%%%%%%%%%%%%%%%%%%%%%%%%%%%%%%%%%%%%%%%%%%%%%%%%%%%%%%%%%%%%
\begin{document}
	

	
\begin{letter} {   
	\hfill Prof. Tullio Vardanega\\
	\hfill Prof. Riccardo Cardin\\
	\hfill Università degli Studi di Padova \\
	\hfill Dipartimento di Matematica ''Tullio Levi Civita'' \\
	\hfill Via Trieste 63 \\
	\hfill 35121 Padova \\
}

\date{}
\opening{Egregio Professore Tullio Vardanega, \newline Egregio Professore Riccardo Cardin,}
\begin{quotation}
	\noindent Con la presente il gruppo \textit{\gruppo} intende comunicarVi ufficialmente che ha intenzione di partecipare alla \emph{Revisione di Accettazione} prevista per il giorno \emph{Venerdì 20 luglio}, al fine di mostrarVi il prodotto finito da Voi commissionato, denominato \textbf{Marvin: dimostratore di Uniweb su Ethereum}, proposto da Red Babel.
	Si allegano alla presente lettera i seguenti documenti:
	\begin{itemize}
		\item \AdR{};
		\item \SdF{};
		\item \G{};
		\item \NdP{};
		\item \PdP{};
		\item \PdQ{};
		\item \MU{};
		\item \MS{};
		\item \VIuno{};
		\item \VEuno{};
		\item \VIdue{};
		\item \VEdue{};
		\item \VItre{};
		\item \VEtre{};
		\item \emph{Verbale\_I\_4\_2018.04.30};
		\item \emph{Verbale\_E\_4\_2018.07.05};
		\item \emph{Verbale\_I\_5\_2018.05.22};
		\item \VIsei{};
	\end{itemize}

	Il codice prodotto è presente al seguente indirizzo:
	\url{https://github.com/SOS-SonsOfSwe/Marvin/}.\\
	
	\newpage
	Inoltre, all’indirizzo internet \url{http://marvinsos.surge.sh} può reperire
una versione già configurata e funzionante del prodotto.
	
	
	Il gruppo Sons of SWE consegna il prodotto richiesto il giorno 19/07/2018, con relativo collaudo, salvo ulteriori accordi, per il giorno 20/07/2018. Il costo finale del prodotto è di \EUR{13.239,00} come descritto all'interno del documento \PdP{}.
	
	Si riportano di seguito i nomi dei componenti del gruppo:\\
	\definecolor{greySOS}{RGB}{209, 222, 223}
	\begin{longtable}{ c  c  c }
		\rowcolor{greySOS}
		\textcolor{white}{\textbf{Nominativo}} & \textcolor{white}{\textbf{Matricola}} & \textcolor{white}{\textbf{Firma}}\\
		%\includegraphics[height=2.5cm]{immagine} aggiustare l'altezza
		Andrea Favero & 1125545 & \includegraphics[height=0.5cm]{firme/AndreaFavero.png} \\
		
		Eleonora Thiella & 1099980 & \includegraphics[height=0.5cm]{firme/EleonoraThiella.png} \\
		
		Federico Caldart & 1097005 & \includegraphics[height=0.5cm]{firme/FedericoCaldart.png} \\
		
		Giovanni Cavallin & 1148957 & \includegraphics[height=0.5cm]{firme/GiovanniCavallin.png} \\
		
		Giovanni Dalla Riva & 1102075 & \includegraphics[height=0.5cm]{firme/GiovanniDallaRiva.png} \\
		
		Lorenzo Menegon & 1097596 & \includegraphics[height=0.5cm]{firme/LorenzoMenegon.png} \\
		
		Stefano Panozzo & 1097068 & \includegraphics[height=0.5cm]{firme/StefanoPanozzo.png} \\
		
	\end{longtable}

Rimaniamo a Vostra completa disposizione per qualsiasi chiarimento in merito ai documenti presentati.

	\end{quotation}
\makeletterclosing{Cordiali saluti,\\}
{
\setlength\parindent{24pt}
\indent i membri del gruppo \gruppo.
}
\end{letter}

\end{document}
	
	
	
	